{\footnotesize
\begin{flushleft}
    La partie « Magnétostatique » aborde l'étude du champ magnétique en régime stationnaire en prenant appui sur les équations locales : la loi de Biot et Savart ne figure pas au programme. L'objectif réside davantage dans l'étude des propriétés des champs magnétiques que dans leur calcul : les calculs de champ magnétique doivent donc se limiter à des situations d'intérêt pratique évident. Enfin, la notion de potentiel-vecteur est hors programme.
\end{flushleft}

    % \begin{minipage}[t]{0.54\textwidth}
        \begin{tabularx}{\linewidth}{|p{0.35\textwidth}|X|}
            \hline 
            \textbf{Notions et contenus} & \textbf{Capacités exigibles} \\
            \hline 
            \multicolumn{2}{|l|}{5.3. Magnétostatique} \\
            \hline 
            Équations locales de la magnétostatique et formes intégrales : flux conservatif et théorème d'Ampère.  & 
            Choisir un contour fermé et une surface et les orienter pour appliquer le théorème d'Ampère. \\
            \hline 
             Linéarité des équations.  & Utiliser une méthode de superposition.  \\
            \hline 
             Propriétés de symétrie. & Exploiter les propriétés de symétrie des sources (rotation, symétrie plane) pour prévoir des propriétés du champ créé. \\
            \hline
            Propriétés topographiques.  & Justifier qu'une carte de lignes de champ puisse ou non être celle d'un champ magnétostatique.
            Repérer, sur une carte de champ magnétostatique, d'éventuelles sources du champ et leur sens.
            Associer l'évolution de la norme d'un champ magnétique à l'évasement des tubes de champ. \\
        \hline 
        \end{tabularx}
% \end{minipage}
% \begin{minipage}[t]{0.43\textwidth}
%     \begin{tabularx}{\linewidth}{|p{0.35\textwidth}|X|}
%         \hline 
        
%         Interactions ion-molécule et molécule-molécule. & Expliquer qualitativement la solvatation des ions dans un solvant polaire. \\
%         \hline 
%         Dipôle induit. Polarisabilité. & Associer la polarisabilité et le volume de l'atome en ordre de grandeur. \\ 
%         \hline
%         Plan infini uniformément chargé en surface. & Établir l'expression du champ créé par un plan infini uniformément chargé en surface.\\ 
%         \hline
%         Établir l'expression du champ créé par un plan infini uniformément chargé en surface. & Établir l'expression du champ créé par un condensateur plan.
%         Déterminer l'expression de la capacité d'un condensateur plan.
%         Citer l'ordre de grandeur du champ disruptif dans l'air.
%         Déterminer l'expression de la densité volumique d'énergie électrostatique dans le cas du condensateur plan à partir de celle de l'énergie du condensateur. \\ 
%         \hline
%         Énergie de constitution d'un noyau atomique modélisé par une boule uniformément chargée. & Énergie de constitution d'un noyau atomique modélisé par une boule uniformément chargée. \\ 
%         \hline
%     \end{tabularx}
% \end{minipage}
\vspace{0.5cm}
}
